% TOP_PROBLEM: {"problemId": "problem.global-regularity-of-the-3d-incompressible-hall-mhd-equations-for-large-data", "canonicalTitle": "Global Regularity of the 3D Incompressible Hall-MHD Equations for Large Data", "releaseRank": 392, "primaryDomain": "analysis_pde", "primaryDomainLabel": "Analysis & PDE", "displayStatus": "open", "releaseStatus": "open"}
% !TeX program = lualatex
% ATTEMPT_STATUS: unresolved
% Model: GPT-6 Astra Ultra; continuation dated 27 September 2026.
\documentclass[11pt]{article}
\usepackage[margin=25mm]{geometry}
\usepackage{fontspec}
\setmainfont{Latin Modern Roman}
\usepackage{amsmath,amssymb,mathtools}
\usepackage{unicode-math}
\setmathfont{Latin Modern Math}
\usepackage{xurl,hyperref}
\hypersetup{colorlinks=true,urlcolor=blue}
\setcounter{secnumdepth}{0}
\setlength{\parindent}{0pt}
\setlength{\parskip}{5pt}
\setlength{\emergencystretch}{3em}
\begin{document}
\section*{\texorpdfstring{Global Regularity of the 3D Incompressible Hall-MHD Equations for Large Data}{Global Regularity of the 3D Incompressible Hall-MHD Equations for Large Data}}\label{392}
\subsection{Definitions and mathematical statement}
A standard isotropic incompressible Hall-MHD system for velocity $u$, magnetic field $b$, and pressure $p$ is
\[
 \begin{aligned}
 \partial_tu+(u\cdot\nabla)u+\nabla p&=\nu\Delta u+(b\cdot\nabla)b,\\
 \partial_tb+(u\cdot\nabla)b-(b\cdot\nabla)u
   +\nabla\times((\nabla\times b)\times b)&=\eta\Delta b,\\
 \nabla\cdot u=\nabla\cdot b&=0.
 \end{aligned}
\]
The Hall term is the curl of the current--field cross product. The question concerns continuing smooth solutions for arbitrary, not necessarily small, smooth initial data for all positive time. \emph{Scope limitation:} the catalogue does not specify the spatial domain, decay or Sobolev conditions, viscosity/resistivity, or isotropic versus partial dissipation. Its reference treats an axisymmetric horizontally dissipative system. The displayed conventional system explains the terminology but does not claim that all these inequivalent choices have been fixed by the supplied record; a precise proof target must retain the chosen source's model and data class.
\par\smallskip\textit{Formulation and concept references:} \href{https://link.springer.com/article/10.1007/s44198-022-00063-8}{[S1]}.

\subsection{Short English statement}
For the intended three-dimensional incompressible Hall-MHD model, do all admissible smooth initial fields remain smooth for all time without a smallness assumption?
\subsection{Research attempt}
\paragraph{A precise conditional scope, without rewriting the catalogue.}
The formulation deliberately leaves several model choices open.
The inspected source [S1] treats horizontal dissipation and a
specified axisymmetric class; its abstract adds restrictions on
the velocity swirl and magnetic vorticity. Its large-data theorem
cannot be cited as a theorem for arbitrary three-dimensional
fields in the displayed isotropic system.

For the calculations below we explicitly choose the displayed
isotropic equations on $\mathbb T^3$, with fixed $\nu,\eta>0$
and smooth divergence-free periodic data. This is an auxiliary
precise model for partial analysis, not a claim that the omitted
catalogue choices have been uniquely recovered. Integrations by
parts also apply to sufficiently decaying smooth fields on
$\mathbb R^3$, but the periodic explicit family below is not
called finite-energy data on $\mathbb R^3$.

We derive the cancellation at energy level, the higher-derivative
Hall remainder, a conditional continuation estimate, and a global
large-data family. The invariant subsystem with zero magnetic field
then identifies an immediate obstruction to a general regularity
proof even in this selected fully dissipative model.

\paragraph{The energy cancellation is exact.}
Write $J=\nabla\times b$. Pair the velocity equation with $u$
and the magnetic equation with $b$, and integrate over the torus.
Divergence-free transport and the pressure contribute zero. The
two magnetic stretching terms cancel, since
\[
 \int u\cdot(b\cdot\nabla)b
       =-\int b\cdot(b\cdot\nabla)u.
\]
For the Hall term, the curl integration identity gives
\[
 \int b\cdot\nabla\times(J\times b)
       =\int J\cdot(J\times b)=0
 \tag{392.1}
\]
pointwise in the final integrand. Thus every smooth solution satisfies
\[
 \frac12\frac{d}{dt}(\|u\|_2^2+\|b\|_2^2)
       +\nu\|\nabla u\|_2^2+\eta\|\nabla b\|_2^2=0.
 \tag{392.2}
\]
The Hall term is nondissipative in this identity. Its disappearance
does not mean it is a lower-order term: it contains two spatial
derivatives and reappears after differentiating the equation.

In a horizontally dissipative model, the same cancellation gives
only horizontal gradient integrals from diffusion. Replacing those
by the full gradients in (392.2) would add control not provided
by the model. This is why the dissipation choice cannot be hidden
in a notation change.

\paragraph{What happens when the magnetic equation is differentiated.}
Let $m\geq3$ be an integer. Apply $D^\alpha$, $|\alpha|\leq m$,
to the magnetic equation and pair with $D^\alpha b$. Moving the
outer curl onto that test function makes the differentiated Hall
contribution
\[
 \int D^\alpha(J\times b)\cdot D^\alpha J.
\]
Its leading term $(D^\alpha J)\times b$ is orthogonal to
$D^\alpha J$, so it cancels exactly. The remainder is
\[
 \int\left[D^\alpha(J\times b)-(D^\alpha J)\times b\right]
                            \cdot D^\alpha J.
 \tag{392.3}
\]
Every product inside the brackets has at least one derivative on
the second factor $b$. A tame product estimate gives
\[
 \|D^\alpha(J\times b)-(D^\alpha J)\times b\|_2
       \leq C_m\|\nabla b\|_\infty\|b\|_{H^m}.
 \tag{392.4}
\]
One can see the derivative count directly for $|\alpha|=m$.
Write the two factors in each Leibniz term as derivatives of
$F=\nabla b$, of respective orders $j$ and $m-1-j$.
The interpolation estimates
\[
 \|D^jF\|_{L^{2(m-1)/j}}
 \leq C\|F\|_\infty^{1-j/(m-1)}
                   \|F\|_{H^{m-1}}^{j/(m-1)}
\]
with the endpoint $j=0$ interpreted as $L^\infty$, followed by
Holder's inequality, bound their product by
$C\|F\|_\infty\|F\|_{H^{m-1}}$. Lower-order terms are handled
the same way with inhomogeneous Sobolev norms. This is the usual
integer tame estimate; no unbounded $m+1$ derivative is hidden
in its right-hand side.

Combining (392.3)--(392.4), summing in $\alpha$ and using Young's
inequality gives
\[
 |\text{Hall contribution}|
 \leq C_m\|\nabla b\|_\infty\|b\|_{H^m}\|\nabla b\|_{H^m}
 \leq\frac\eta2\|\nabla b\|_{H^m}^2
    +\frac{C_m}{\eta}\|\nabla b\|_\infty^2\|b\|_{H^m}^2.
 \tag{392.5}
\]
The square of the Lipschitz norm, and the reciprocal resistivity,
are material. An estimate with only the energy norm of $b$ in this
coefficient would be a stronger assertion not proved by the cancellation.

\paragraph{A conditional higher-regularity bound.}
Put $Y_m=\|u\|_{H^m}^2+\|b\|_{H^m}^2$. For the ordinary
transport and stretching terms, the highest differentiated
transport cancels against itself by divergence freedom, while the
two highest $b$-coupling terms cancel against each other exactly
as at energy level. The remaining Leibniz terms are bounded by
the same tame interpolation mechanism, now with only one derivative
in the coefficients. Together with (392.5), this gives
\[
 \begin{split}
 \frac{dY_m}{dt}
   +\nu\|\nabla u\|_{H^m}^2+\eta\|\nabla b\|_{H^m}^2
 \leq C_{m,\nu,\eta}
  \bigl(\|\nabla u\|_\infty+\|\nabla b\|_\infty
                      +\|\nabla b\|_\infty^2\bigr)Y_m.
 \end{split}
 \tag{392.6}
\]
For clarity, a typical transport commutator has the form
$D^\alpha(u\cdot\nabla v)-u\cdot\nabla D^\alpha v$.
Its pairing is bounded by
$C(\|\nabla u\|_\infty\|v\|_{H^m}
+\|\nabla v\|_\infty\|u\|_{H^m})\|v\|_{H^m}$;
summing the velocity and magnetic versions yields the linear
Lipschitz coefficients in (392.6). The pressure disappears under
the divergence-free differentiated velocity test.

Gronwall's inequality proves that $Y_m$ stays bounded up to any
finite time $T$ for which
\[
 \int_0^T\bigl(\|\nabla u\|_\infty+
               \|\nabla b\|_\infty+\|\nabla b\|_\infty^2\bigr)dt<\infty.
 \tag{392.7}
\]
Together with a local existence/continuation theorem in the chosen
Sobolev class, this is a conditional continuation criterion. The
energy identity does not establish (392.7); this is the first
missing a priori bound in this route.

Indeed the Sobolev embedding $H^m\hookrightarrow W^{1,\infty}$
in three dimensions holds for $m>5/2$. Substituting that estimate
into (392.6) yields a differential upper bound containing $Y_m^2$.
Such an ODE comparison gives local control but permits finite-time
blowup of its comparison solution. It is not a global regularity
argument merely because the original fields also obey (392.2).

\paragraph{The energy space does not supply a Lipschitz bound.}
Choose a nonzero smooth compactly supported divergence-free field
$B$ in a small Euclidean ball; for example take the curl of a
compactly supported smooth vector potential. Embed its small rescalings
inside a coordinate patch of the torus and set
\[
 B_\lambda(x)=\lambda^{1/2}B(\lambda(x-x_0)).
\]
For large $\lambda$ the support fits in the patch, and ordinary
Euclidean scaling gives
\[
 \|B_\lambda\|_2^2=\lambda^{-2}\|B\|_2^2,\qquad
 \|\nabla B_\lambda\|_2^2=\|\nabla B\|_2^2,
 \quad
 \|\nabla B_\lambda\|_\infty
      =\lambda^{3/2}\|\nabla B\|_\infty.
 \tag{392.8}
\]
Thus even simultaneous boundedness of the instantaneous $L^2$
and $H^1$ norms would not bound the required Lipschitz coefficient.
The actual energy law only integrates the $H^1$ seminorm in time,
which is weaker still. This scaling example is not a singular
solution; it rejects a proposed functional-inequality shortcut.

\paragraph{An exact global large-data magnetic family.}
On the periodic domain let
\[
 u=0,\qquad b(x,t)=(a(x_3,t),c(x_3,t),0),
\]
where $a,c$ solve the one-dimensional heat equation with diffusivity
$\eta$ and arbitrary smooth periodic initial values. Divergence
freedom is immediate. Since $b\cdot\nabla=a\partial_1+c\partial_2$,
we have $(b\cdot\nabla)b=0$, so the velocity equation holds with
constant pressure. Moreover
\[
 J=(-\partial_3c,\partial_3a,0),\qquad
 J\times b=(0,0,-a\partial_3a-c\partial_3c).
\]
The last field is a gradient in the $x_3$ direction, whose curl
vanishes. Hence the Hall term is exactly zero, and the magnetic
equation reduces to the stated heat equations.

Fourier series give
$a(x_3,t)=\sum_{\ell\in\mathbb Z}a_\ell e^{-\eta\ell^2t}e^{i\ell x_3}$,
and similarly for $c$. Smooth initial coefficients decay faster
than every power, so these series remain smooth for all $t\geq0$.
There is no smallness requirement on their amplitude. This is an
actual global large-data family of the displayed isotropic periodic
equations, but it does not describe generic three-dimensional
magnetic fields or the source's separate symmetry class.

\paragraph{The zero-magnetic-field subsystem is an exact obstruction.}
If $b\equiv0$, the magnetic equation is identically satisfied and
the velocity equation is precisely the three-dimensional
incompressible Navier--Stokes equation
\[
 \partial_tu+(u\cdot\nabla)u+\nabla p=\nu\Delta u,
 \qquad\nabla\cdot u=0.
 \tag{392.9}
\]
Thus a theorem proving global smoothness for arbitrary smooth
divergence-free data of the fully dissipative Hall--MHD model would
in particular prove global smoothness for this velocity subsystem.
The Hall cancellation cannot remove that difficulty, since it
vanishes identically on this subsystem. This observation is a
reduction, not a claim that Hall--MHD and Navier--Stokes are
equivalent in all directions.

One might try to use a nonzero magnetic field to regularize the
velocity and then pass to zero field. Such a passage would require
bounds uniform as the magnetic amplitude tends to zero. Without
them the argument cannot address (392.9). A regularity result for
selected magnetic geometry or large stabilizing background fields
therefore does not answer the arbitrary-data question either.

\paragraph{Diagnostics and outcome.}
The accompanying symbolic checks verify the pointwise cross-product
cancellation, the exact one-coordinate magnetic ansatz, and the
scaling exponents in (392.8). They do not simulate a blowup or
prove that (392.7) holds for general solutions.

The established partial results are the full energy identity, the
explicit higher-derivative Hall estimate and conditional criterion,
and a genuinely large-data global family in a precisely declared
isotropic periodic model. The first unresolved analytic step is a
global bound replacing the uncontrolled coefficients in (392.7).
At the catalogue level, a final theorem must also specify and
retain the intended domain, dissipation and data class. Neither
the present auxiliary calculations nor the distinct source theorem
settles every choice left open by the printed record.
\subsection{Sources}
\begin{itemize}
\item[S1]Li and Cui, Global Well-Posedness for the 3D Axisymmetric Hall-MHD System with Horizontal Dissipation, 2022, Introduction.
\url{https://link.springer.com/article/10.1007/s44198-022-00063-8}.
\textit{Formulation source.}
\end{itemize}
\end{document}
